\section{Introduction}

The value of a summary is rarely intrinsic; it is realized only when the summary helps a particular reader accomplish a particular goal. Yet automatic summarization has long been evaluated through metrics that aim to capture a general notion of quality as a property of the summary itself. Lexical overlap measures like ROUGE compare candidate summaries against reference texts~\cite{lin2004rouge}, embedding-based approaches like BERTScore assess semantic similarity~\cite{zhang2020bertscore}, and more targeted metrics evaluate specific dimensions such as factual consistency, coherence, or readability~\cite{kryscinski2019neural, wang2020asking, fabbri2021summeval, luo2022readability}. More recently, Large Language Models (LLMs) have been deployed as judges, promising more nuanced assessments that approach human-level evaluation~\cite{zheng2023judging,liu2023geval}. Across this proliferation of metrics, however, quality is treated as a fixed attribute that can be measured against a gold standard or set of reference criteria, independent of who is reading.

This framing obscures a fundamental truth about summarization: a summary is only useful insofar as it serves the informational needs of the person reading it. Consider a recent paper on mRNA vaccine development. A biomedical researcher may need a summary that preserves methodological detail, statistical results, and connections to prior literature, while a family physician reading the same paper may be best served by a summary emphasizing clinical implications, patient guidance, and contraindications. The source document is the same in each case, and a reference summary written for one audience may score highly under standard metrics while failing entirely to satisfy another reader's needs. Existing evaluation paradigms have little to say about this mismatch, because they do not model the reader at all.

To capture this reader-centered notion of utility, we introduce \textit{information satisfaction} as an axis of summarization evaluation. Information satisfaction measures the extent to which a summary resolves the specific informational need that motivated the reader to consult the source in the first place, operationalized in this chapter as a query paired with a persona describing the reader's role and expertise. Unlike generic quality metrics, which ask whether a summary is fluent, faithful, or similar to a reference, information satisfaction asks if a summary meets a specific user's informational needs.
A summary can be fluent, factually accurate, and a faithful condensation of its source while still leaving the reader's question unanswered, and conversely a stylistically rough summary may fully satisfy the query at hand for a reader with the right background. By foregrounding both the query and the reader behind it, information satisfaction reframes evaluation around the reader's purpose rather than the summary's surface properties.

Given this reframing, we ask whether existing summarization metrics are equipped to measure information satisfaction at all. We organize our investigation around two research questions:

\paragraph{{\chsevenrqone\label{chsevenrq1}} Are evaluation metrics sensitive to variation in informational content and reader persona?}
A metric that meaningfully tracks information satisfaction should respond when the information in a summary shifts relative to what the reader needs—either because the summary's content changes or because the persona consuming it changes. We probe this property by systematically varying both axes and measuring how popular automatic metrics, including traditional reference-based measures and LLM-as-judge approaches, respond to the perturbations.

\paragraph{{\chsevenrqtwo\label{chsevenrq2}} Do evaluation metrics agree with human judgments of information satisfaction?} 
Even a metric that is sensitive to content and persona changes is only useful if its judgments align with those of real readers. We conduct an expert human evaluation in which annotators rate summaries with respect to a specified query and persona, and we measure the correlation between these human judgments and the scores produced by automatic metrics.


In this chapter, we show that many popular metrics, including strong LLM-as-judge metrics, fail simple perturbation tests to the informational content of a summary. Given this finding, we unsurprisingly find that the same metrics show poor agreement with human judges of information satisfaction. Taken together, our results suggest that neither traditional nor LLM-based metrics adequately capture information satisfaction, and that progress on user-centered summarization will require evaluation frameworks that explicitly model who the summary is for and what they came to learn.
